\section{Descripción del caso de uso}

Se eligió un sistema RAG para consultar documentación técnica y se lo adaptó a
una distribuidora mayorista. La información está repartida entre documentos y
datos operativos. Encontrar una respuesta puede llevar tiempo, depender de
alguien que conozca el proceso o terminar en el uso de una versión que ya no
está vigente.

Se propone la capa de datos que permitiría a un copiloto interno recuperar
información y usarla como fuente para responder. El alcance incluye
almacenamiento, integridad, autorización, recuperación y trazabilidad. El
entrenamiento de modelos y la aplicación completa quedan fuera de la prueba.

Se consideran tres perfiles. \textbf{Operaciones y Logística} consulta fichas de
producto, procedimientos e incidencias autorizadas. \textbf{Comercial y
Compras} consulta condiciones comerciales, políticas de venta e información de
proveedores. \textbf{Administración y Calidad} accede a documentación legal y
de cumplimiento, y revisa la trazabilidad. Los servicios de ingesta, consulta y
administración documental no son usuarios finales y reciben los permisos
mínimos para su tarea.

La base registra versiones documentales, identifica cuál está vigente, divide
su contenido en fragmentos y asocia embeddings. También resuelve consultas
estructuradas sobre pedidos, entregas o condiciones comerciales. Cada respuesta
exitosa guarda la evidencia que identifica la fuente utilizada.

Los riesgos principales son recuperar contenido no autorizado, usar una versión
sustituida o revocada y responder sin una fuente verificable. Para evitarlos, se
conserva el historial, se aplican los permisos antes de la recuperación y se
guarda evidencia y auditoría de cada interacción.

El alcance se limita a la capa de datos. No se incluyen frontend, backend, API,
autenticación, integración con un LLM ni embeddings generados por un modelo
real. Esos componentes podrían consumir la base en una etapa futura.
